% ==============================
% 60_service_orchestration.tex
% Kubernetes Cluster integration
% ==============================

\chapter{Containerized Service Deployment}\label{chap:deployment}
\section{Basic K3s Fundamentals}
While \texttt{containerd} handles the low-level execution of containers, \gls{k3s} manages their orchestration across the entire cluster. 

Kubernetes utilizes \texttt{pods} as the smallest deployable units. In a cluster environment, pods are ephemeral by design, which introduces the "Challenges of Ephemerality" that must be addressed for a stable environment:

\begin{itemize}
	\item \textbf{Dynamic Networking} \\
	Pods receive internal IP addresses randomly upon creation. To ensure reliable communication and reverse proxying, we utilize \textbf{Services} (ClusterIP), which provide a persistent DNS name regardless of individual pod restarts.
	
	\item \textbf{Data Persistence} \\
	By default, data within a pod's volume is lost if the pod shifts to a different node[cite: 115, 116]. This is resolved in Section \ref{section:pod-storage} through the use of Persistent Volume Claims (PVCs).
	
	\item \textbf{Database Integrity} \\
	Managing SQL databases in a distributed cluster is complex. While standard files can live on network shares, databases require specific locking mechanisms to prevent corruption[cite: 113].
\end{itemize}


\section{K3s Features}
By default, K3s comes with many features that you may or may not use.

\begin{itemize}
	\item \textbf{Ingress Controller} \\
	It comes with a modified version of \texttt{Traefik} that handles \texttt{INGRESS} traffic. (to expand upon)
	
	\item \textbf{Integrations} \\
	In K3s, you can add \texttt{operators}, which are essentially plugins for Kubernetes that integrates with other services.
\end{itemize}

\section{Using K3s Manifest}
In \texttt{Kubernetes}, you declare a manifest file that is used to define the ideal state of the cluster. Kubernetes will then work towards that goal by orchestrating your services accordingly.

This makes the cluster being self-healing and declarative a core function of the server.

\section{Package Manager}
\texttt{K3s} comes with \texttt{Helm}, a package manager for Kubernetes.

\section{Storage Orchestration}\label{section:pod-storage}
Please read chapter \ref{chap:storage} on storage if you haven't already. Many core concepts visited there will be used here.

Depending on your choice of storage (Distributed or Centralized), there are different ways to achieve working services without losing data whenever a pod changes node.

If you use centralized storage :
\begin{itemize}
	\item \textbf{Network File Shares (NFS/SMB)} \\
	Ideal for configuration files and media libraries. 
	By mounting a centralized \gls{zfs} dataset (e.g., \texttt{/mnt/appdata}) to every node, pods can read and write to the same physical location regardless of their host.
	
	\item \textbf{SQL Servers} \\
	Many services require or support SQL servers, which can be hosted on the central storage server. Often, services support a range of versions, so if you're lucky, you can squeeze a few in the same server. 
	
	For exemple, \texttt{Vaultwarden} uses SQLite but can support MariaDB as well. 
	
	\item \textbf{The SQLite/SQL Challenge} \\
	SQLite databases are prone to corruption over NFS due to improper file-locking.
	For these workloads, we utilize \textbf{Democratic-CSI} to provision block-storage (\texttt{zvol}) via \texttt{iSCSI}. 
	This provides the pod with a dedicated virtual disk that supports the low-level locking required for database stability.
\end{itemize}



\section{Network Orchestration}\label{section:pod-network}
This section covers the networking configuration for the cluster as well as the access to the services deployed. The latter is the focus because internal networking doesn't require much configuration.

\begin{itemize}
	\item \textbf{Internal Traffic} () \\
	Services within the cluster can communicate between themselves using DNS names, much like docker containers. 
	There is no need to keep track of IPs or to meddle with that.
	
	\item \textbf{External Traffic (Ingress)} \\
	To have access to your services, you must go through \texttt{K3s} built-in \texttt{Ingress} controller, built on a modified version of \texttt{Traefik}. This makes sure requests are forwarded to the right pod and is the layer that the \gls{reverse proxy} interacts with.
\end{itemize}




\subsection{Secure Remote Access}
It is reccommended to set up secure remote access to reach 'admin' services that should never be opened to the public, such as your server's dashboard. 

\begin{itemize}
	\item \textbf{VPN Access (Zero-Exposure / Zero-Tier)} \\
	Tools like Tailscale shine here if you do not need or want to push your services publicly. Things like admin services in particular.
	\item \textbf{Reverse Proxy (Public-Facing)} \\
	If you do want to push your services to the public, a reverse proxy is the way to go. It allows you to expose only your designated services and doesn't even touch the rest.
	\item \textbf{Hybrid-Theory (VPN Access + Reverse Proxy)} \\
	You can always combine both. It works very well. If your reverse proxy does not have direct access to your ingress controller, you can connect to it from the tailscale IP (because they are technically on the same network). Otherwise, having it in case your proxy goes down for any reason is good too.
\end{itemize}